You have 50 minutes to finish this part. Use only the pencil to make drawings
and mark-ups on the sky map. After finishing, fill your student ID on the
answer sheet and the sky map. Put the two papers in the folder, and leave
the compass, ruler, and pencil on the table.

Use the sky map found in the envelope. The map is in azimuthal projection,
so distance from the centre is proportional to zenith distance. The
instructions for marking points, constellations, and curves are shown below.
\ioaafig{2022_OBS-SKY_Q13-f1.pdf}
\textit{The separately issued blind sky map and the referenced IAU-code
Table 1 are absent from the supplied question PDF. Annotated charts survive
in the official solution and are reproduced there.}

The map represents the sky somewhere in Romania. The zenith is at its centre.
\begin{enumerate}
\item Find the latitude of the place and draw the greatest circumpolar circle
  on the map (Mark A).
\item Draw with continuous lines the celestial equator (Mark B), the ecliptic
  (Mark C), the galactic equator (Mark G), and the local meridian (Mark M).
\item Mark the cardinal points (N, S, E, W) and show three Solar System
  planets (P1, P2, P3).
\item Identify and mark the five brightest stars, with magnitudes less than
  $+1.00$, numbered from 1 (brightest) to 5 (faintest). Give their Bayer names.
\item Draw ten constellations you can identify, labelled with the corresponding
  IAU codes from Table 1.
\item Mark M31 (number 6), M27 (7), M15 (8), $\beta$ Cet (9), and $\delta$ UMi (10).
\item Estimate the sidereal time of the map.
\item Estimate Altair's ($\alpha$ Aquilae) equatorial coordinates.
\end{enumerate}
